\section{Related Work}
\label{sec:related}

\paragraph{DID identification and its threats.}
A large econometrics literature formalizes when DID estimators fail under
staggered timing and heterogeneous effects: the ``forbidden'' comparisons of
two-way fixed effects~\cite{goodmanbacon2021}, negative
weights~\cite{dechaisemartin2020}, robust
estimators~\cite{callaway2021,sunabraham2021}, assumption and pre-trend
surveys~\cite{roth2023}, and fragile inference under serial
correlation~\cite{bertrand2004trust}. \sys{} adds no estimator: it
\emph{operationalizes} this literature as an auditable rubric, turning each
known threat into a dimension whose supporting evidence can be checked in a
given paper (\S\ref{sec:framework}).

\paragraph{LLM causal reasoning and benchmarks.}
A parallel line asks whether LLMs reason about causality. Benchmarks
(CLadder~\cite{jin2023cladder}, Corr2Cause~\cite{jin2024corr2cause},
CausalBench~\cite{zhou2024causalbench}) and
analyses~\cite{kiciman2023causal,yang2024criticalreview,zecevic2023causalparrots}
report mixed evidence: strong performance on some causal tasks, but models
often reproduce causal \emph{language} without reliable causal
\emph{inference}; a few move toward applied settings and toward separating
\emph{identification} from
\emph{estimation}~\cite{lee2025econcausal,sawarni2026causalreasoning}. \sys{}
asks the inverse question: are a \emph{human} paper's identification assumptions
adequately evidenced? A causal verdict is out of scope by design.

\paragraph{Responsible AI for empirical research.}
AI is increasingly applied to research reliability: reproducibility checking,
statistical-reporting validation, and AI-assisted review, motivated by wide
analyst-driven result spreads~\cite{silberzahn2018many,botviniknezer2020variability}
and specification searching in causal economics~\cite{brodeur2020methods}. LLMs
are also used directly as evaluators and reviewers, judging model
outputs~\cite{zheng2023judge}, giving feedback on research
papers~\cite{liang2024feedback}, and reviewing inside automated-science
pipelines~\cite{lu2024aiscientist}; these systems produce open-ended critique,
whereas \sys{} scores a fixed rubric against in-paper evidence, so its outputs
are comparable across papers and measurable against injected ground truth. \sys{}
differs in target and posture: it audits \emph{identification credibility}
specifically (the part of an argument that licenses a causal claim) and issues
no verdicts, producing evidence-grounded, human-reviewable reports that localize
risk for the expert's final call. Closest to our setting,
CausalVerify~\cite{causalverify} grounds evaluation in \emph{execution}: it runs
a model's generated estimation code against synthetic data and checks the
recovered effect. That benchmark places identification-assumption auditing
outside its scope, which is the layer \sys{} takes up; we reuse its released
corpus of published economics papers and add a dimension-level
identification-risk annotation layer on top of it.

\paragraph{NLP for climate evidence.}
Climate-NLP systems verify claims against evidence
(CLIMATE-FEVER~\cite{climatefever2020}), detect environmental
claims~\cite{stammbach2023environmental}, analyse corporate climate disclosures
with traceable LLM answers~\cite{ni2023chatreport}, and tune models for
faithful evidence-based question answering~\cite{schimanski2024faithful}.
These systems ground answers in documents; \sys{} extends document-grounded
verification from \emph{claims} to the \emph{identification argument} of a
causal study, auditing whether reported evidence supports each assumption.

\paragraph{Evaluation methodology and bounded agents.}
Our evaluation follows benchmarking lessons: holistic
evaluation~\cite{liang2023helm} over a single number (we report a layered
retrieval/relevance/adequacy/risk profile); dynamic, adversarial
benchmarking~\cite{kiela2021dynabench} (our model-in-the-loop flaw injection);
and the caution that operationalizing a capability is
non-neutral~\cite{raji2021benchmark,gebru2021datasheets}
(\S\ref{sec:discussion}). Where ReAct-style agents~\cite{yao2023react} act over
an open horizon, \sys{} \emph{bounds} the paradigm (agency confined to two
scoped sub-tasks behind deterministic control flow), trading autonomy for
reproducibility and evidence-traceability.
